\subsection{Problem Formalization}


\label{sec:prelim:formalization}
The objects fixed above assemble into one problem statement. Let $G$ be an AIG with the
node roles of~\eqref{eq:prelim:role} and edge inversions of~\eqref{eq:prelim:inversion} as
its features, and let $S$ be a deterministic synthesis algorithm
(\ref{sec:prelim:algorithms:orchestrate}) mapping $G$ to the function-equivalent,
optimized $S(G)$. The prediction target is the \emph{optimizability} of $G$ under $S$,
the fraction of nodes the algorithm removes,
\begin{equation}
    y(G) = 1 - \frac{\left| \mathcal{V}(S(G)) \right|}{\left| \mathcal{V}(G) \right|}
    \;\in\; [0, 1].
    \label{eq:prelim:target}
\end{equation}
Obtaining $y(G)$ means running $S$, which is the expensive step prediction is meant to
avoid. An encoder $f_\theta$, a GNN in the sense of~\ref{sec:prelim:gnn:mp} with a
graph-level readout, is trained to approximate $y$. A reduction operator $R$, drawn from
the three families of~\ref{sec:prelim:reduction}, shrinks the encoder's input.

RQ5 (\ref{sec:intro:rq5}) is well posed only under a compatibility requirement: $R(G)$ and
$G$ must occupy the same input space, so that one set of weights $\theta$ ingests both.
Every summarization method here therefore produces features of which the full-graph
features are the single-member special case (\ref{sec:method:summary:schema}). A
reduction that changed the schema would make cross-state inference undefined rather than
merely inaccurate.